% Exact staging-only proof for production Rethlas verification.
% Date: 29 August 2026.
% This file is not a live-book input and makes no novelty or priority claim.

\begin{proof}[Proof of
\cref{prop:p21p3-rigid-exception-extension}]
Because $X$ is separated, its underlying topological space is Hausdorff.
A locally finite union of closed subsets of a Hausdorff space is closed;
applied to the singletons in $E$, this shows that $E$ is closed.  Hence
$Y=X\setminus E$ is an open analytic subspace of $X$.  Goodness,
separatedness, and strictness pass to open analytic subspaces.  An open
immersion has empty relative boundary, and the boundary-composition formula
\cite[(11.4.1.3)]{CLD25} therefore shows that $Y$ is boundaryless.
We shall also use the corresponding separation property for currents.
Every good separated analytic space is locally holomorphically separated by
\cite[\S15.2.2]{CLD25}, and is therefore locally differentially separated by
\cite[Corollary~15.2.4]{CLD25}.  This applies to $X$ and to its open analytic
subspaces, in particular to $Y$ and to the open sets considered below.

We next prove the assertion about connectedness.  By
\cite[Theorem~3.5.1]{Ducros24CurveStructure}, the underlying topological
space of every analytic curve is a graph.  A connected graph is connected by
injective arcs by \cite[\S1.3.5]{Ducros24CurveStructure}.  On the other hand,
\cite[Proposition~3.5.2]{Ducros24CurveStructure} states that the number of
branches at a rigid point of an analytic curve equals the cardinality of its
inverse image in the normalization.  Since $X$ is normal, every point of $E$
is consequently unibranch.  By
\cite[\S1.3.12]{Ducros24CurveStructure}, every point lying in the interior of
an injective arc is pluribranch.

The set $Y$ is nonempty.  Indeed, if $E$ is empty this is immediate; if
$x\in E$, the fact that $x$ is unibranch gives a nonempty branch issuing
from $x$.  Choose a connected component of $U\setminus\{x\}$ representing
that branch for a sufficiently small neighborhood $U$ of $x$, and choose a
point $z$ in that component.  After shrinking $U$, the graph structure
supplies an injective arc from $x$ to $z$ with nonempty interior.  Any point
in that interior is pluribranch, hence does not belong to $E$.  Now let
$y_0,y_1\in Y$.  If $y_0=y_1$, there is nothing to prove.
Otherwise choose an injective arc $I$ in $X$ joining them.  If $I$ met $E$,
any point of $I\cap E$ would be an interior point of $I$, because its
endpoints lie in $Y$.  Such a point would be both pluribranch and unibranch,
a contradiction.  Thus $I\subset Y$.
This proves \cref{itm:p21p3-rigid-exception:connected}, in fact by showing
that $Y$ is connected by injective arcs.

Let $T$ be as in
\cref{itm:p21p3-rigid-exception:classification}.  Restriction to the open
subspace $Y$ is defined by extension by zero of compactly supported test
forms, as in \cite[\S7.1.6]{CLD25}.  It preserves positivity; by the locality
and differential identities of \cite[\S\S7.7.1--7.7.3]{CLD25}, it also
preserves $\mathrm d'$- and $\mathrm d''$-closedness.  The curve $Y$ has all
the hypotheses of the already established rig-smooth classification
\cref{prop:p21p3-gfc-smooth-rigidity}.  Hence there is a unique
$c\in\mathbf R_{\geq0}$ such that
\begin{equation}
\label{eq:p21p3-rigid-exception-on-complement}
                            T|_Y=c\delta_Y.
\end{equation}
By the definition of the integration current in
\cite[\S7.3.1]{CLD25} and the restriction operation of
\cite[\S7.1.6]{CLD25}, one has $\delta_X|_Y=\delta_Y$.  Put
$S:=T-c\delta_X$.  Then $S|_Y=0$ by
\eqref{eq:p21p3-rigid-exception-on-complement}.  Let $\alpha$ be an
arbitrary compactly supported smooth $(1,1)$-test form on $X$.  A rigid
point $x$ satisfies $d_K(x)=0$: indeed, rigid means that
$\mathscr H(x)/K$ is finite by \cite[\S11.1.10]{CLD25}, and both the residue
transcendence-degree term and the rational-rank term in the dimension formula
\cite[(11.2.3.1)]{CLD25} then vanish.  This remains true for a finite
inseparable extension.  The analytic form-support statement
\cite[Proposition~15.2.1(a)]{CLD25}, which is the analytic consequence of
the tropical support vanishing in \cite[Proposition~3.3.6]{CLD25}, gives
\begin{equation}
\label{eq:p21p3-rigid-exception-test-support}
 \operatorname{Supp}(\alpha)
 \subseteq\{x\in X:d_K(x)\geq1\}
 \subseteq X\setminus E=Y.
\end{equation}
Thus $\alpha|_Y$ is a compactly supported test form on $Y$, and therefore
\[
                   \langle S,\alpha\rangle
                   =\langle S|_Y,\alpha|_Y\rangle=0.
\]
Since $\alpha$ was arbitrary, $S=0$.  This proves
\eqref{eq:p21p3-rigid-exception-classification}.  Uniqueness of $c$ follows
after restriction to $Y$ from uniqueness in
\cref{prop:p21p3-gfc-smooth-rigidity}.

Finally, let $W,u,v$ satisfy
\eqref{eq:p21p3-rigid-exception-equal-curvature}, and put $h=u-v$.
The local differential separation recorded in the first paragraph puts us
within the hypotheses of the order-zero result.  Positive currents have
order zero by
\cite[Proposition~7.2.9]{CLD25}, so multiplication by the continuous
function $h$ is defined as in \cite[Example~7.2.12]{CLD25}.  Corollary~7.2.10
of \cite{CLD25} defines $[h]$ for the locally integrable function $h$.  For
every compactly supported smooth $(1,1)$-test form $\alpha$ on $W$, the
definitions in \cite[Corollary~7.2.10, Example~7.2.12, and
\S7.3.1]{CLD25} give
\[
 \bigl\langle h(\delta_X|_W),\alpha\bigr\rangle
 =\int_W h\alpha
 =\langle[h],\alpha\rangle.
\]
Consequently $h(\delta_X|_W)=[h]$.  Using
\eqref{eq:p21p3-rigid-exception-classification} and linearity gives
\begin{align*}
 \mathrm d'\mathrm d''\bigl(hT|_W\bigr)
 &=c\,\mathrm d'\mathrm d''[h]\\
 &=c\bigl(\mathrm d'\mathrm d''[u]
          -\mathrm d'\mathrm d''[v]\bigr)=0.
\end{align*}
This is \eqref{eq:p21p3-rigid-exception-bt-conclusion} and proves
\cref{itm:p21p3-rigid-exception:bedford-taylor}.
\end{proof}
